\documentclass[10pt,a4paper]{amsart}
\usepackage[margin=19mm]{geometry}
\usepackage{amsmath,amssymb,mathtools}
\usepackage[scr=rsfs,cal=euler]{mathalfa}
\usepackage{xfrac}
\usepackage[unicode,hidelinks,bookmarks=false]{hyperref}
\newcommand{\ie}{i.e.\ }
\newcommand{\cf}{cf.\ }
\newcommand{\resp}{respectively\ }
\newcommand{\Subsection}{\subsection}
\providecommand{\nobreakdash}{\nobreak\hbox{-}}
\setlength{\emergencystretch}{2em}
\multlinegap0pt
\usepackage[all]{xy}
\newtheorem{lmm}{Lemma}
\newtheorem{thm}{Theorem}
\def\Mfld{{\rm Mfld}}
\def\Perf{{\rm Perf}}
\def\cE{\mathcal E}\def\cF{\mathcal F}\def\cG{\mathcal G}\def\cH{\mathcal H}
\def\fr{{\rm fr}}
\def\ot{\otimes}
\newcommand{\ra}{\rightarrow}
\def\tr{{\rm tr}}
\newcommand{\xra}{\xrightarrow}
\title{Chern--Simons factorization algebras and knot polynomials}
\author{Kevin Costello, John Francis, Owen Gwilliam}
\date{Source: arXiv:2602.12412v1 (CC BY 4.0)}
\begin{document}
\maketitle

\noindent\emph{Source and licence.} Chern--Simons factorization algebras and knot polynomials. Version arXiv:2602.12412v1 (CC BY 4.0), distributed by arXiv under the Creative Commons Attribution 4.0 International licence, http://creativecommons.org/licenses/by/4.0/, as recorded in the arXiv licence metadata for this exact version and shown on the abstract page https://arxiv.org/abs/2602.12412v1. Full licence notice: \texttt{licenses/arxiv-2602.12412-CC-BY-4.0.txt}.\par\smallskip

\noindent\textit{Editorial scope.} Counted unit: the article's thm thm (source0.tex lines 1137--1148). The article's complete original proof of it is delivered with it (source0.tex lines 1150--1156, one \texttt{proof} environment of 7 lines). No mathematics is written, repaired or re-derived: the statement and the proof are the article's own lines, in the article's own notation, and no retained line is substituted or re-flowed.  Retained uncounted as the source-local context the proof needs: lines 1099--1106.  Nothing was omitted from any retained span, so the package carries no omission record.  No retained line contains a citation, so the wrapper carries no bibliography and no bibliography entry.  No retained line contains a figure environment, an image-inclusion command or drawing code, so no image is delivered and the package declares no asset dependency.  Editorial formatting only, in addition: an AMS \texttt{amsart} preamble that restates the article's own declarations and macros used by the retained lines, namely the environments \texttt{lmm}, \texttt{thm} on separate counters, together with the standard packages those lines need.  The retained environments print in this document as Lmm 1, Theorem 1 in order of appearance, whereas the article prints the same items with the numbering of its own full document.  The complete native source of the article as distributed in the arXiv e-print for this exact version is bundled unchanged as \texttt{sources/194547/source0.tex}.
\begin{lmm}\label{lemma.bimmodule}
Let $(A,R)$ and $(A,S)$ be $\cE_{1\subset n}$-algebras, 
and let $V$ be an $(A;R,S)$-bimodule that is perfect as an $S$-module. There is a natural map
\[
\int_{K\subset M}(A,R) \to \int_{K\subset M}(A,S)
\]
of factorization homology theories on $\Mfld_{1\subset n}^{\fr}$ after restriction to {\bf closed} 1-manifolds~$K$.
\end{lmm}

\begin{thm}  
Let $M$ be a framed $n$-manifold with $K\subset M$ a framed link. 
Let $(A,R)$ be an $\cE_{1\subset n}$-algebra with $V$ an $(A;R, A)$-bimodule that is perfect as an $A$-module. 
There is a natural map ${\sf m}_V:\int_{K\subset M}(A,R)\to \int_MA$ such that the diagram 
\[
\xymatrix{
&\displaystyle\int_K R\ar[rr]\ar[dd]&&\displaystyle\int_{K\subset M}(A,R)\ar[dd]^-{{\sf m}_V}\\
\Bbbk\ar[rd]_{\tr(V)}\ar[ru]^-{\rm unit}\\
&\displaystyle\int_{K}A\ar[rr]&& \displaystyle\int_{M}A\\}
\]
commutes up to equivalence.
\end{thm}

\begin{proof} 
Applying Lemma~\ref{lemma.bimmodule} to the case of $S=A$ gives the commutative square above. It remains to establish the commutative triangle on the left. By definition the map $\int_K R \ra \int_K A$ is given, for each connected component of $K$, by the value of Hochschild homology on the map
\[
\Perf_R \xra{-\ot_R V} \Perf_A~.
\]
This functor sends the unit $R\in\Perf_R$ to $V$. Applying Hochschild homology, this sends the unit element of $\int_{S^1} R$ to the element $\tr(V)\in \int_{S^1}A$.
\end{proof}

\end{document}
